\documentclass[10pt,a4paper]{amsart}
\usepackage[margin=19mm]{geometry}
\usepackage{amsmath,amssymb,mathtools}
\usepackage[scr=rsfs,cal=euler]{mathalfa}
\usepackage{xfrac}
\usepackage[unicode,hidelinks,bookmarks=false]{hyperref}
\newcommand{\ie}{i.e.\ }
\newcommand{\cf}{cf.\ }
\newcommand{\resp}{respectively\ }
\newcommand{\Subsection}{\subsection}
\providecommand{\nobreakdash}{\nobreak\hbox{-}}
\setlength{\emergencystretch}{2em}
\multlinegap0pt
\usepackage{bm}
\usepackage{bbm}
\usepackage{dsfont}
\usepackage{xspace}
\usepackage{cancel}
\usepackage{mathtools}
\usepackage{amsthm}
\makeatletter
\@ifundefined{theorem}{\newtheorem{theorem}{Theorem}}{}
\@ifundefined{lemma}{\newtheorem{lemma}[theorem]{Lemma}}{}
\makeatother
\newcommand{\eigen}{\operatorname{eigen}}
\newcommand{\nw}{\mathrm{new}}
\title[Effective Hecke eigenvalue equidistribution over the Atkin--]{Effective Hecke eigenvalue equidistribution over the Atkin--Lehner subspaces}
\author{Aarya J. Kumar, Sargam Mondal, Erick Ross, Hui Xue}
\date{Source: arXiv:2609.17806v1 (CC BY 4.0)}
\begin{document}
\maketitle

\noindent\emph{Source and licence.} Effective Hecke eigenvalue equidistribution over the Atkin--Lehner subspaces. Version arXiv:2609.17806v1 (CC BY 4.0), distributed under the Creative Commons Attribution 4.0 International licence, as recorded in the arXiv licence metadata for this exact version and shown on the abstract page https://arxiv.org/abs/2609.17806v1. Full rights notice: licenses/arxiv-2609.17806-CC-BY-4.0.txt.\par\smallskip

\noindent\textit{Editorial scope.} Counted unit: the Lemma whose source label is \texttt{lem:sum-upper-bound} in the article (source0.tex lines 577--583, source label \texttt{lem:sum-upper-bound}), delivered together with the article's own complete original proof of it (source0.tex lines 585--666, one \texttt{proof} environment of 82 lines). No mathematics is written, repaired or re-derived: statement and proof are the article's own text line for line. Required context retained uncounted: lem:trace (lines 422--435); lem:weyl-limits (lines 492--498). Every definition, display and label of those lines is retained unchanged. Omitted from this package and not redistributed: 1--421, 436--491, 499--576, 584--584, 667--2472. Nothing inside a retained excerpt is omitted or altered: every retained body is delivered exactly as the article's lines are written, so no omission receipt of any kind accompanies this package. Citations: the retained excerpts carry no citation command at all, so no bibliography entry of the article is transcribed and no reference is redistributed. Reference adaptation: none was needed and none was made; every reference inside the delivered unit resolves inside it, so no unresolved reference or citation is produced. Printed slips kept literal: no printed word, symbol, hypothesis or number of the counted statement or of its proof was altered. Layout and class: the article's own class is \texttt{amsart} with packages including fullpage, amsmath, amsthm, amssymb, amscd, enumerate, enumitem, array, float, bbm, bm, stmaryrd, comment, mathtools; the wrapper is the lane's pinned \texttt{amsart} editorial wrapper at 10pt on A4, which restates the theorem-like environments the retained text needs and adds the article's own macro definitions that the retained text needs (\texttt{\textbackslash eigen}, \texttt{\textbackslash nw}), with extra wrapper packages (bm, bbm, dsfont, xspace, cancel, mathtools). The delivered numbering is the unit's own. Assets: no retained span contains a figure environment, a graphics-inclusion command, a PDF-page inclusion, a table or a TikZ picture; no image file is copied into this package, the package declares no asset dependency and no image file is redistributed with it. Licence: version arXiv:2609.17806v1 of this article is distributed under the Creative Commons Attribution 4.0 International licence, as recorded in the arXiv licence metadata for this exact version and shown on the abstract page https://arxiv.org/abs/2609.17806v1. A full rights notice is delivered at licenses/arxiv-2609.17806-CC-BY-4.0.txt. Characters: every retained excerpt is pure ASCII, so the delivered engine, XeLaTeX with its default Latin Modern text font, reports no missing character.
\begin{lemma}\label{lem:trace}
    For coprime $m, N \in \mathbb{N}$ and even weight $k \geq 2$,
    $$\left| \operatorname{Tr}_{S_k^\sigma(N)}\mathbf{T}_m' - M(k, m, N) \right| \leq E(m, N)$$
    and
    $$\left| \operatorname{Tr}_{S_k^{\nw, \sigma}(N)}\mathbf{T}_m' - M^{\nw}(k, m, N) \right| \leq E(m, N).$$
Here,
\begin{align*}
    M(k, m, N) &= \frac{\mathds{1}_{m = \square}}{\sqrt{m}} \cdot \frac{k - 1}{12} \cdot \frac{\psi(N)}{2^{\omega(N)}}\\
    M^{\nw}(k, m, N) &= \frac{\mathds{1}_{m = \square}}{\sqrt{m}} \cdot \frac{k - 1}{12} \cdot \frac{\psi^{\nw}(N)}{2^{\omega(N)}} \prod_{p^r || N} \left( 1 + \sigma(p^r) \cdot \frac{-\mathds{1}_{r = 2}}{p^2 - p - 1}  \right)\\
    E(m, N) &= 6.7261717 \cdot  m \sigma_0(m) \sigma_0(N)^2 \log(N + 2) \sqrt{N}\\
    \psi(N) &= N \prod_{p \mid N} \left(1 + \frac{1}{p} \right)\\
    \psi^{\nw}(N) &= N \prod_{p \mid N} \left(1 - \frac{1}{p} -  \frac{\mathds{1}_{r \geq 2}}{p^2} + \frac{\mathds{1}_{r \geq 3}}{p^3} \right).
\end{align*}
\end{lemma}

\begin{lemma} \label{lem:weyl-limits}
    For any $m \in \mathbb{Z}$, we have
    $$ \left| \sum_{\lambda \in \eigen_{S_k^\sigma(N)}(\mathbf{T}_p')} 2\cos(m\theta_\lambda) -c_m r^\sigma \right| \leq  2E(p^{|m|},N) + 
    \mathds{1}_{2 \mid m} \cdot \left| p^{-|m|/2} - p^{-(|m| - 2)/2} \right| E(1, N).
    $$
    The same result also holds over $S_k^{\nw, \sigma}(N).$
\end{lemma}

\begin{lemma}\label{lem:sum-upper-bound}
    For all $M \in \mathbb{N}$ and sub-intervals $I \subseteq [-2, 2],$ we have
    \begin{align}
    &\quad \sum_{1 \leq |m| \leq M} \left( \frac{1}{M + 1} + \frac{1}{\pi|m|} \right) \left| \sum_{\lambda \in \eigen_{S_k^\sigma(N)}(\mathbf{T}'_p)}  2\cos\left(m \theta_\lambda \right) - r^\sigma c_m \right| \label{eqn:sum-from-lemma-sum-upper-bound}\\
    &\le 70.9375 \cdot p^M \cdot N^{1/2} \sigma_0(N)^2\log(N + 2).
    \end{align}
\end{lemma}

\begin{proof}
    Since the cosine and the Weyl limit are even functions in $m,$ we can rewrite \eqref{eqn:sum-from-lemma-sum-upper-bound} as
\[
\eqref{eqn:sum-from-lemma-sum-upper-bound} =
2\sum_{1\leq m\leq M}
\left(
\frac{1}{M+1}
+
\frac{1}{\pi m}
\right)
\left|
\sum_{\lambda \in \eigen_{S_k^\sigma(N)}(\mathbf{T}'_p)} 2\cos(m\theta_{\lambda})-c_m r^\sigma
\right|.
\]

By Lemmas \ref{lem:trace} and \ref{lem:weyl-limits},
\begin{align*}
    & \quad \left|
\sum_{\lambda \in \eigen_{S_k^\sigma(N)}(\mathbf{T}'_p)} 2\cos(m \theta_\lambda )-c_m r^\sigma
\right| \\
&\leq 2\cdot 6.7261717 \cdot  p^m \sigma_0(p^m) \sigma_0(N)^2 \log(N + 2) \sqrt{N} \\
&\qquad + \mathds{1}_{2 \mid m} \cdot 6.7261717 \,  \sigma_0(N)^2 \log(N + 2) \sqrt{N} \cdot \left| p^{-m/2} - p^{-(m - 2)/2}  \right|  \\
&= \left( (m + 1)p^m  +  \frac{\mathds{1}_{2 \mid m}}{2}\left| p^{-m/2} - p^{-(m - 2)/2}  \right|\right) \cdot 2 \cdot 6.7261717  \,  \sigma_0(N)^2 \log(N + 2) \sqrt{N}.
\end{align*}
Therefore, we obtain
\begin{align}
    \eqref{eqn:sum-from-lemma-sum-upper-bound} &\leq 26.904687 \,\sigma_0(N)^2 \log(N + 2) \sqrt{N}\\
    & \quad\ \times \left[ \sum_{1 \leq m \leq M} \left( \frac{1}{M + 1} + \frac{1}{\pi m} \right) (m + 1)p^m \right. \label{eqn:sum-upper-bound-expression-1}\\
    & \qquad\quad + \left. \sum_{\substack{0 < m \leq M\\ \text{$m$ even}}} \left( \frac{1}{M + 1} + \frac{1}{\pi m} \right) \cdot \frac{1}{2}\left| p^{-m/2} - p^{-(m - 2)/2} \right| \right] . \label{eqn:sum-upper-bound-expression-2}
\end{align}

First, we bound \eqref{eqn:sum-upper-bound-expression-1}. Observe that
\begin{align}
    & \eqref{eqn:sum-upper-bound-expression-1} \\
    &= \frac{1}{M + 1} \left(  \sum_{1 \leq m \leq M} p^m + \sum_{1 \leq m \leq M} mp^m\right) + \frac{1}{\pi} \left( \sum_{1 \leq m \leq M} p^m  + \sum_{1 \leq m \leq M} \frac{p^m}{m} \right)\\
    &= \frac{1}{M + 1} \left( \frac{p(p^M - 1)}{p -1} + \frac{M p^{M+ 2} - (M + 1)p^{M + 1} + p}{(p - 1)^2} \right) + \frac{1}{\pi} \left(\frac{p(p^M - 1)}{p -1}  + \sum_{1 \leq m \leq M} \frac{p^m}{m} \right)\\
    &=  \left(  \frac{(M + 1) p^{M+ 2} - (M + 2)p^{M + 1} + 2p - p^2}{(M + 1)(p - 1)^2} \right) + \frac{1}{\pi} \left( \frac{p^{M + 1}}{p - 1}  + \sum_{1 \leq m \leq M} \frac{p^m}{m} \right) - \frac{1}{\pi} \cdot \frac{p}{p - 1}\\
    &=  \left(  \frac{(M + 1) \left( p^{M + 2} - p^{M + 1} \right)}{(M + 1)(p - 1)^2} - \frac{p^{M + 1}}{(M + 1)} \right) + \frac{1}{\pi} \left(\frac{p^{M + 1}}{p - 1} + \sum_{1 \leq m \leq M} \frac{p^m}{m} \right) - \left( \frac{1}{\pi} + \frac{p^2 - 2p}{(M + 1)(p - 1)^2} \right)\\
    &=  \left( 1 + \frac{1}{\pi} \right) \frac{ p^{M + 1} }{p - 1} - \frac{p^{M + 1}}{(M + 1)}  + \frac{1}{\pi}  \sum_{1 \leq m \leq M} \frac{p^m}{m} - \left( \frac{1}{\pi} + \frac{p^2 - 2p}{(M + 1)(p - 1)^2} \right). \label{eqn:simplified---sum-upper-bound-expression-1}
\end{align}
Additionally, note that
\begin{align*}
    \frac{ \sum_{1 \leq m \leq M} \frac{p^m}{m}}{\frac{p^{M + 1}}{(M + 1)}} &= (M + 1) \sum_{1 \leq m \leq M} \frac{p^{m - M - 1}}{m} \\
    &\leq \frac{1}{p} \sum_{k = 1}^{M - 1} \frac{M + 1}{M - k} \cdot  \frac{1}{p^k} \qquad  \text{substituting $k = M - m$} \\
    &\leq \frac{1}{2} \sum_{k = 1}^{M - 1} \frac{M + 1}{M - k} \cdot  \frac{1}{2^k} \qquad \text{since $p \geq 2$} \\
    &\leq \frac{1}{2} \left( \sum_{k = 0}^{M - 2} \frac{k + 3}{2} \cdot \frac{1}{2^k}  +  \frac{M + 1}{2^{M - 1}} \right)\\
    &< \frac{1}{2} \left( \sum_{k = 0}^{\infty} \frac{k + 3}{2} \cdot \frac{1}{2^k} +  2 \right) \\
    &< \frac{1}{2} \left( \frac{1}{2}\sum_{k = 0}^{\infty} k \cdot \frac{1}{2^k} + \frac{3}{2} \sum_{k = 0}^{\infty} \frac{1}{2^k} +  2 \right) \\
    &\leq \frac{1}{2} \left( 1 + 3 + 2 \right) \qquad  \text{since $\sum_{k = 0}^\infty k x^k = \frac{x}{(1 - x)^2}$}\\
    &< \pi.
\end{align*}
This implies that $- \frac{p^{M + 1}}{(M + 1)}  + \frac{1}{\pi}  \sum_{1 \leq m \leq M} \frac{p^m}{m} < 0,$ so we have from \eqref{eqn:simplified---sum-upper-bound-expression-1} that
\begin{align}
     \eqref{eqn:sum-upper-bound-expression-1} &\leq \left( 1 + \frac{1}{\pi} \right)\frac{ p^{M + 1} }{p - 1} - \left( \frac{1}{\pi} + \frac{p^2 - 2p}{(M + 1)(p - 1)^2} \right) \\
    &\leq \left(2 + \frac{2}{\pi} \right) p^M - \left( \frac{1}{\pi} + \frac{p^2 - 2p}{(M + 1)(p - 1)^2} \right) \label{eqn:simplified-expression-2}.
\end{align}

Second, we bound \eqref{eqn:sum-upper-bound-expression-2}. Observe that for all $M \in \mathbb{N},$
\begin{align}
    \eqref{eqn:sum-upper-bound-expression-2} &\leq \sum_{\substack{0 < m \leq M\\ \text{$m$ even}}} \left( \frac{1}{M + 1} + \frac{1}{\pi} \right) \cdot \frac{1}{2}\left( p^{-(m - 2)/2} - p^{-m/2}\right)\\
    &= \frac{1}{2}\left( \frac{1}{M + 1} + \frac{1}{\pi} \right) \left( 1 - \frac{1}{p^{\lfloor M/2 \rfloor}} \right)\\
    &\leq  \frac{1}{\pi} + \frac{p^2 - 2p}{(M + 1)(p - 1)^2}.
\end{align}
The final inequality is verified through the following cases.
\begin{itemize}
    \item If $M \geq 3,$ then $M + 1 > \pi,$ so
    $$\frac{1}{2}\left( \frac{1}{M + 1} + \frac{1}{\pi} \right) \left( 1 - \frac{1}{p^{\lfloor M/2 \rfloor}} \right) \leq \frac{1}{\pi} \left( 1 - \frac{1}{p^{\lfloor M/2 \rfloor}} \right) \leq \frac{1}{\pi} \leq \frac{1}{\pi} + \frac{p^2 - 2p}{(M + 1)(p - 1)^2}.$$
    \item If $M = 1,$ then $\frac{1}{2}\left( \frac{1}{M + 1} + \frac{1}{\pi} \right) \left( 1 - \frac{1}{p^{\lfloor M/2 \rfloor}} \right) = 0.$
    \item If $M = 2$ and $p \leq 43,$ then
    $$\frac{1}{2}\left( \frac{1}{M + 1} + \frac{1}{\pi} \right) \left( 1 - \frac{1}{p^{\lfloor M/2 \rfloor}} \right) \leq \frac{1}{\pi} \leq \frac{1}{\pi} + \frac{p^2 - 2p}{(M + 1)(p - 1)^2}.$$
    \item If $M = 2$ and $p \geq 47,$ then
    $$\frac{1}{2}\left( \frac{1}{M + 1} + \frac{1}{\pi} \right) \left( 1 - \frac{1}{p^{\lfloor M/2 \rfloor}} \right) \leq \frac{1}{2} \left( \frac{1}{3} + \frac{1}{\pi} \right) \leq \frac{1}{\pi} + \frac{47^2 - 2 \cdot 47}{3 (47 - 1)^2} \leq \frac{1}{\pi} + \frac{p^2 - 2 \cdot p}{(M + 1) (p - 1)^2},$$
    since $\frac{1}{\pi} + \frac{p^2 - 2 \cdot p}{(M + 1) (p - 1)^2}$ is increasing in $p$ for $M = 2.$
\end{itemize}

Combining the above bounds for \eqref{eqn:sum-upper-bound-expression-1} and \eqref{eqn:sum-upper-bound-expression-2}, we obtain $\eqref{eqn:sum-upper-bound-expression-1} + \eqref{eqn:sum-upper-bound-expression-2} \leq \left( 2 + \frac{2}{\pi} \right) p^M,$ so that
\begin{align}
    \eqref{eqn:sum-from-lemma-sum-upper-bound} &\le 26.904687 \, \sigma_0(N)^2 \log(N + 2) \sqrt{N} \Big[\eqref{eqn:sum-upper-bound-expression-1} + \eqref{eqn:sum-upper-bound-expression-2}\Big] \\
    &\le 70.9375 \cdot p^M \cdot \sigma_0(N)^2 \log(N + 2) \sqrt{N},
\end{align}
as desired.
\end{proof}

\end{document}
